
%\needspace{4\baselineskip}
\color{uniblue}{\subsection{Дистанционная защита (ДЗ)}}
\color{black}

\begin{enumerate}[label=\arabic{section}.\arabic{subsection}.\arabic{enumi}, labelsep=4pt, leftmargin=0pt, itemindent=57pt, itemsep=0pt, parsep=5pt] % если что можно крутить itemsep
%%%%%%%%%%%%%%%%%%%%%%%%%%%%%%%%%%%%%%%%%%%%%%%%%%%%%%%%%%%%%%%%%%%%%%  ОБЩИЕ СВЕДЕНИЯ %%%%%%%%%%%%%%%%%%%%%%%%%%%%%%%%%%%%%%%%%%%%%%%%%%
\item Общие сведения

\begin{enumerate}[label=\arabic{section}.\arabic{subsection}.\arabic{enumi}.\arabic*, labelsep=4pt, leftmargin=0em, itemindent=65pt, parsep=0pt]

\item
Дистанционная  защита  (ДЗ)  предназначена  для  селективного  отключения  поврежденного  элемента электрической  сети  при  всех видах  коротких  замыканий (КЗ),  реагируя  на снижение измеряемого сопротивления в защищаемой зоне.

\item 
ДЗ селективно срабатывает:
\begin{itemize}
    \item при всех видах КЗ, в том числе при постановке линии под напряжение; 
    \item при переходе внешнего КЗ во внутреннее, в том числе при реверсе мощности; 
    \item при КЗ, сопровождающемся снижением измеряемого напряжения до нуля.
\end{itemize}

\item 
Обеспечивается несрабатывание защиты:
\begin{itemize} 
\item при качаниях и асинхронном ходе; 
\item при неисправностях цепей напряжения;
\item при внешних КЗ с насыщением трансформаторов тока (при соблюдении заявленных производителем требований к ТТ; 
\item при постановке линии под напряжение и включении в транзит в отсутствие КЗ;
\item при реверсе мощности при внешних КЗ; 
\item при бросках тока намагничивания без КЗ;
\item при внешних КЗ, даже если измеряемое напряжение снижается до нуля.
\end{itemize}

\item Каждая из ступеней ДЗ от междуфазных КЗ включает в себя три реле сопротивления (РС), использующие для замера сопротивления  комплексные значения линейных напряжений и соответствующих им линейных токов:
% ФОРМУЛЫ
\begin{align}
\vec Z_{AB} =  \frac{\vec U_A - \vec U_B}{\vec I_A - \vec I_B}\label{eq:groupZ1}, \\[10pt]
\vec Z_{BC} =  \frac{\vec U_B - \vec U_C}{\vec I_B - \vec I_C}\label{eq:groupZ2}, \\[10pt]
\vec Z_{CA} =  \frac{\vec U_C - \vec U_A}{\vec I_C - \vec I_A}\label{eq:groupZ3}, 
\end{align}
\begin{eqexpl}[25mm]
    \item{$ \vec U_{A}, \vec U_{B}, \vec U_{C} $} комплексные значения напряжений фаз А, В, С;
    \vspace{2pt}
    \item{$ \vec I_{A}, \vec I_{B}, \vec I_{C} $} комплексные значения токов фаз А, В, С.
\end{eqexpl}
\vspace{4pt}

\item 
Комплексные значения сопротивлений по контуру «фаза-земля» рассчитываются с учетом компенсации тока нулевой последовательности своей линии и параллельной линии:
% ФОРМУЛЫ
\begin{align}
\vec Z_{A} =  \frac{\vec U_A }{\vec I_A + \vec K \cdot  3 \vec I_0 + \vec K_{\textsl{ПЛ}} \cdot 3 \vec I_{\textsl{0ПЛ}}}\label{eq:groupZ11}, \\[10pt]
\vec Z_{B} =  \frac{\vec U_B }{\vec I_B + \vec K \cdot  3 \vec I_0 + \vec K_{\textsl{ПЛ}} \cdot 3 \vec I_{\textsl{0ПЛ}}}\label{eq:groupZ12}, \\[10pt]
\vec Z_{C} =  \frac{\vec U_C }{\vec I_C + \vec K \cdot  3 \vec I_0 + \vec K_{\textsl{ПЛ}} \cdot 3 \vec I_{\textsl{0ПЛ}}}\label{eq:groupZ13}, 
\end{align}
\begin{eqexpl}[25mm]
    \item{$ 3 \vec I_{0} $}  комплексное  значение  утроенного  тока  нулевой  последовательности  защищаемой  линии (расчетное значение);
    \item{$ 3 \vec I_{\textsl{0ПЛ}} $}   комплексное  значение  утроенного  тока  нулевой  последовательности  параллельной  линии (измеренное значение);
    \item{$ \vec K $}  комплексный коэффициент компенсации тока нулевой последовательности защищаемой линии;
    \item{$ \vec K_{\textsl{ПЛ}} $}    комплексный  коэффициент  компенсации  тока  нулевой  последовательности  параллельной линии.
\end{eqexpl}
\vspace{4pt}

\item 
В устройстве коэффициент  $\vec K$  задается двумя уставками, \textbf{«KR - коэффициент компенсации 3I0 своей линии»}  и  \textbf{«KХ  -  коэффициент  компенсации  3I0  своей  линии»}. коэффициент $ \vec K_{\textsl{ПЛ}}$   задается  уставками \textbf{«KR// - коэффициент компенсации 3I0 параллельной линии»} и \textbf{«KХ// - коэффициент компенсации 3I0 параллельной линии»}. 

\item 
Расчет  коэффициентов  компенсации  зависит  от  режима  работы  параллельной  линии.  Если параллельная линия в работе, отключена или отсутствует, то коэффициенты компенсации рассчитываются по следующим формулам:
% ФОРМУЛЫ
\begin{align}
\vec K = \frac{1}{3} \cdot \frac{\vec Z_{0 \textsl{уд}} - \vec Z_{1 \textsl{уд}}}{\vec Z_{1 \textsl{уд}}}\label{eq:groupK1}, \\[10pt]
\vec K_{\textsl{ПЛ}} = \frac{1}{3} \cdot \frac{\vec Z_{M \textsl{уд}}}{\vec Z_{1 \textsl{уд}}}\label{eq:groupK2}, 
\end{align}
\begin{eqexpl}[25mm]
\item{$ \vec Z_{0 \textsl{уд}} $}  комплексное удельное сопротивление нулевой последовательности;
\item{$ \vec Z_{1 \textsl{уд}} $}  комплексное удельное сопротивление прямой последовательности;
\item{$ \vec Z_{M \textsl{уд}} $}  комплексное удельное сопротивление взаимоиндукции.
\end{eqexpl}
\vspace{4pt}

\item 
Компенсация  влияния  тока  параллельной  линии  блокируется,  если  величина  тока  нулевой последовательности  параллельной  линии $ 3 \vec I_{\textsl{0ПЛ}} $  превышает  135\%  от  величины  тока  нулевой последовательности защищаемой линии $ 3 \vec I_{0} $:
% ФОРМУЛА
\begin{equation}
3 \vec I_{\textsl{0ПЛ}} >1,35 \cdot 3 \vec I_{0},
\end{equation}

\item 
Если параллельная линия отключена и заземлена, то коэффициент  $ \vec K $  рассчитывается по формуле:
% ФОРМУЛА
\begin{equation}
\vec K = \frac{1}{3} \cdot \left( \frac{\vec Z_{0\textsl{уд}} - \vec Z_{1\textsl{уд}}}{ \vec Z_{1\textsl{уд}}}- \frac{\left(\frac{\vec Z_{M\textsl{уд}}}{\vec Z_{1\textsl{уд}}}\right)^2}{\frac{\vec Z_{0\textsl{уд}} - \vec Z_{1\textsl{уд}}}{\vec Z_{1\textsl{уд}}}+1} \right),
\end{equation}
а коэффициент  $\vec K_{\textsl{ПЛ}}$  в данном случае равен 0.

\item 
Характеристики  срабатывания  ненаправленного  РС  для  междуфазных  и  фазных  контуров  задаются независимо для каждой ступени уставками:
\begin{itemize} 
\item \textbf{«Rср ДЗ-ФФ(ФЗ)-N»} – активное сопротивление срабатывания РС; 
\item \textbf{«Хср ДЗ-ФФ(ФЗ)-N»} – реактивное сопротивление срабатывания РС;
\item \textbf{«Ф1 ДЗ-ФФ(ФЗ)-N»} – угол линии электропередач; 
\end{itemize}
\begin{eqexpl}[3mm]
    \item{$ N $}  номер ступени ДЗ.
\end{eqexpl} 

\item 
Для  предотвращения  неселективного  срабатывания  РС-ФЗ-1  и  РС-ФФ-1  при  внешних  повреждениях через  переходное  сопротивление  предусмотрена  возможность  отсечения  области  характеристики срабатывания,  задаваемой  углом  $\phi_4$.  Область  отсекается  отрезком,  исходящей  под  углом  $\phi_4$,  из  точки  на верхней части характеристики, соответствующей углу линии электропередачи  $\phi_1$. Угол $\phi_4$ задается уставкой \textbf{«Ф4 ДЗ-ФФ(ФЗ)-1»}. 
\item 
Вид характеристики срабатывания ненаправленного РС представлен на рисунке \ref{pict_dz:char1}. 

\begin{figure}[H]
\centering
\includegraphics[width=0.4\textwidth,height=0.4\textheight,keepaspectratio]{\fbpath/561.0021 ДЗ/img/_char1.pdf}
\captionof{figure}{ Характеристика срабатывания ненаправленного РС}\label{pict_dz:char1}
\end{figure}

\item 
Параметры реле сопротивления приведены в таблице \ref{desc_dz:tbl1}

\small
\setlength{\extrarowheight}{0.05cm} % высота ячеек в таблице
\begin{longtable}
{|>{\centering\arraybackslash}m{15.2cm}|>{\centering\arraybackslash}m{1.7cm}|}
\caption{Параметры РС\hfill\vspace{-0.5\baselineskip}}\label{desc_dz:tbl1}\\ 
\hhline{|-|-|}
\rowcolor{gray!30}
Наименование параметра & Значение  \\ 
\hhline{|-|-|}
\endfirsthead

\caption*{\hspace{3pt}\emph{Продолжение таблицы \ref{desc_dz:tbl1}\hfill\vspace{-0.5\baselineskip}}} \\ % сделано по ГОСТ 2.105 п.6.8.7
\hhline{|-|-|}
\rowcolor{gray!30}
ГНаименование параметра & Значение  \\ 
\hhline{|-|-|}
\endhead
\hhline{|-|-|}

\endfoot
\endlastfoot

\raggedright Минимальное  междуфазное  напряжение,  при  котором  обеспечиваются  точностные характеристики РС, В, не менее & \centering\arraybackslash 0,5  \\  \hline
\raggedright Ток точной работы во всем диапазоне уставок при работе на угле φ 1 , \% от номинального тока, не более  & \centering\arraybackslash 10  \\  \hline
\raggedright Собственное  время  срабатывания  РС  при  работе  на  угле  максимальной  чувствительности, трехкратном токе точной работы и скачкообразном уменьшении напряжения  на входе РС от напряжения,  соответствующего  сопротивлению  на  зажимах  РС  не  менее  $1,2\cdot(X_{\textsl{УСТ}} /sin\phi_1 )$ до напряжения соответствующего $0,6\cdot(X_{\textsl{УСТ}} /sin\phi_1 )$ мс, не более   & \centering\arraybackslash 25  \\  \hline
\raggedright Собственное  время  возврата  РС  при  работе  на  угле  максимальной  чувствительности, трехкратном  токе  точной  работы  и  скачкообразном  увеличении  напряжения  на  входе  РС  от напряжения соответствующего сопротивлению на зажимах РС $0,1\cdot(X_{\textsl{УСТ}} /sin\phi_1 )$ до напряжения соответствующего $1,2\cdot(X_{\textsl{УСТ}} /sin\phi_1 )$ (но не более 100 В для междуфазного канала и 57 В –- для фазного), мс, не более   & \centering\arraybackslash 50  \\  \hline
\raggedright Средняя основная погрешность всех РС по величине сопротивления срабатывания $R_{\textsl{УСТ}}$  и $X_{\textsl{УСТ}}$ при токе, равном $I_{\textsl{ном}}$ (или, в зависимости от уставки, меньшем токе, исходя из максимального напряжения на зажимах РС, равного 100 В), \% от уставки, не более  & \centering\arraybackslash $\pm$5  \\  \hline
\raggedright Основная абсолютная погрешность РС по углу наклона характеристики при токе, равном $I_{\textsl{ном}}$, эл. град., не более   & \centering\arraybackslash $\pm$5  \\  \hline
\raggedright Дополнительная абсолютная погрешность РС по углу наклона характеристики при изменении тока КЗ в диапазоне $(0,2 \dots 30) \cdot I_{\textsl{ном}}$, эл. град. от значения, измеренного при $I_{\textsl{ном}}$, не более   & \centering\arraybackslash $\pm$7  \\  \hline
\raggedright Дополнительная  погрешность  РС  по  сопротивлениям  $R$  и  $X$  при  изменении  температуры окружающей  среды  в  рабочем  диапазоне,  \%  от  среднего  значения,  измеренного  при нормальных климатических условиях, не более   & \centering\arraybackslash $\pm$3  \\  \hline
\raggedright Время  срабатывания  ДЗ  при  переходе  внешнего  КЗ  во  внутреннее  в  условиях  наличия насыщения ТТ мс, не более & \centering\arraybackslash 60  \\  \hline
\raggedright Коэффициент возврата, не более   & \centering\arraybackslash 1,1  \\  \hline

\end{longtable}
\setlength{\extrarowheight}{0.0cm} % возврат параметра высота ячеек в таблице
\normalsize

\item 
Выполнение функций направленного РС обеспечивается совместным использованием ненаправленного РС и органов направленности (ОН). Орган направленности представляет собой два отрезка, исходящие из начала координат и расположенные во втором и четвёртом квадрантах. Вид характеристики  (см. рис. \ref{pict_dz:on}) определяется углами $\phi_2$ и $\phi_3$  наклона  этих  отрезков,  отсчитываемых  относительно  оси  $R$.  Углы  наклона  характеристики задаются уставками \textbf{«Ф2 ОН-ФФ(ФЗ)»} и\textbf{«Ф3 ОН-ФФ(ФЗ)»}. Уставки органов направленности являются общими для всех ступеней, задаются раздельно для контуров «фаза-фаза» и «фаза-земля». Углы наклона в устройстве задаются  для  прямонаправленного  ОН,  а  характеристика  обратнонаправленного  ОН  симметрична характеристике прямонаправленного ОН с поворотом на 180 градусов.

\begin{figure}[H]
\centering
\includegraphics[width=0.7\textwidth,height=0.7\textheight,keepaspectratio]{\fbpath/561.0021 ДЗ/img/_on.pdf}
\captionof{figure}{Характеристики срабатывания органов направления}\label{pict_dz:on}
\end{figure}

\item
Направленность для каждой ступени ДЗ задается с помощью программных переключателей \textbf{«Направл. ДЗ-ФФ(ФЗ)-N»}, имеющих три положения:
\begin{enumerate}[label=\arabic*)]
\item Ненапр. – ступень выполнена ненаправленной; 
\item Прямо – ступень выполнена прямонаправленной; 
\item Обратно – ступень выполнена обратнонаправленной. 
\end{enumerate}

\item
При  близких  трехфазных  КЗ,  когда  все  напряжения  на  входе  РС  близки  к  нулю,  для  определения направленности в течение времени не менее 40 мс используются напряжения предаварийного режима (работа по «памяти»). При работе РС «по памяти» при трехфазных КЗ в месте установки защиты обеспечивается длительность сигнала срабатывания на выходе РС не менее 40 мс в диапазоне токов от $0,2\cdot I_{\textsl{ном}}$ до $30\cdot I_{\textsl{ном}}$. При КЗ «за спиной» с токами до $30\cdot I_{\textsl{ном}}$  обеспечивается отсутствие ложных срабатываний РС, в том числе при снижении измеряемого напряжения до нуля. 

\item
Для отстройки от нагрузочного режима возможен вырез в характеристике срабатывания РС. Уставки измерительного органа отстройки от нагрузки являются общими для всех ступеней, задаются раздельно для контуров  «фаза-фаза»  и  «фаза-земля».  Его  характеристика  срабатывания  приведена  на  рисунке \ref{pict_dz:r}.  Эта характеристика срабатывания определяется уставками: 
\begin{itemize}
    \item \textbf{«Rнг ДЗ-ФФ(ФЗ)»} - сопротивление выреза нагрузочного режима; 
    \item \textbf{«Фнг ДЗ-ФФ(ФЗ)»} - угол выреза нагрузочного режима.     
\end{itemize}

\item
Исключаемая область в характеристике срабатывания симметрична относительно осей координат. Если нет  необходимости  отстройки  от  нагрузки,  то  \textbf{«Rнг  ДЗ-ФФ(ФЗ)»}  выбирается  больше  уставки  по  $R$  самой чувствительной ступени ДЗ. 

\begin{figure}[H]
\centering
\includegraphics[width=0.6\textwidth,height=0.6\textheight,keepaspectratio]{\fbpath/561.0021 ДЗ/img/_r.pdf}
\captionof{figure}{Характеристика срабатывания измерительного органа отстройки от нагрузки}\label{pict_dz:r}
\end{figure}

\item
Суммарная характеристика направленного РС, используемого в логике работы ступеней ДЗ приведена на рисунке \ref{pict_dz:rs}.  

\begin{figure}[H]
\centering
\includegraphics[width=0.85\textwidth,height=0.85\textheight,keepaspectratio]{\fbpath/561.0021 ДЗ/img/_rs.pdf}
\captionof{figure}{Характеристики срабатывания направленного РС}\label{pict_dz:rs}
\end{figure}

\item
Дополнительно для каждой ступени ДЗ предусмотрены направленные РС с охватом начала координат, уставки  по  $R$,  $X$  и $\phi_1$  которых  совпадают  с  уставками  ненаправленного  РС  соответствующей  ступени. Характеристика срабатывания (см. рис \ref{pict_dz:nachalo}) направленного РС с охватом начала координат имеет форму параллелограмма, смещенного вдоль направления металлического КЗ на линии в третий и четвертый квадранты на величину до $0,1\cdot X_{\textsl{ср}}$.  

\begin{figure}[H]
\centering
\includegraphics[width=0.4\textwidth,height=0.4\textheight,keepaspectratio]{\fbpath/561.0021 ДЗ/img/_nachalo.pdf}
\captionof{figure}{Характеристика срабатывания направленного РС с охватом начала координат}\label{pict_dz:nachalo}
\end{figure}

\item
Для  обнаружения  однофазных  КЗ  и  пуска  ступеней  ДЗ-ФЗ,  контролирующих  контуры  «фаза-земля», предусмотрен  орган  определения  вида  повреждения  (ООВП),  основанный  на  совместном  применении измерительного органа тока нулевой последовательности (ИО 3I0 ООВП) с торможением от одного из фазных токов (средний из  $I_A , I_B , I_C$) и измерительного органа напряжения нулевой последовательности (ИО 3U0 ООВП). Торможение осуществляется от тока той фазы, значение в которой является средним между максимальным и минимальным значениями токов в остальных двух фазах. В таком случае при междуфазном КЗ на землю торможение максимальном, а при однофазном КЗ –- минимально. 
\item
Характеристика  срабатывания  ИО  3I0  ООВП  с  торможением  представлена  на  рисунке  \ref{pict_dz:torm}.  Порог срабатывания горизонтального участка задается уставкой \textbf{«3I0ср ООВП»}. Наклонный участок характеристики проходит через точку, равной $1,25 \cdot I_{\textsl{ном}}$ , под углом к горизонтальной оси, соответствующим коэффициенту торможения $k_T$. Коэффициент торможения задается уставкой \textbf{«Кт ООВП»}. При превышении тормозного тока уставки \textbf{«Iср блок. торм.»} ООВП блокируется.

\begin{figure}[H]
\centering
\includegraphics[width=0.65\textwidth,height=0.65\textheight,keepaspectratio]{\fbpath/561.0021 ДЗ/img/_torm.pdf}
\captionof{figure}{Характеристика срабатывания ИО 3I0 ООВП}\label{pict_dz:torm}
\end{figure}

\item
Ввод  в  работу  ступеней  ДЗ,  работающих  по  контуру  «фаза-земля»  может  осуществляться  с дополнительным  контролем  напряжения  нулевой  последовательности  за  счет  ввода  программного  ключа \textbf{«Контр.  3U0»}.  Напряжение  срабатывания  ИО  3U0  ООВП  задается  уставкой  \textbf{«3U0ср  ООВП»}. ИО 3U0 ООВП может работать по напряжению $3U_{\textsl{0 звезда}}$, рассчитанному на основе фазных напряжений, или по напряжению $3U_{\textsl{0 треуг.}}$   --  выбирается  программным  ключом  \textbf{«Выбор  3U0»},  имеющим  два  соответствующих положения: \textsl{«Расчетн.»} и \textsl{«Измерен.»}.

\item
Переменная $3U_{\textsl{0 треуг.}}$ зависит от положения программного переключателя \textbf{«Схема подключения цепей \textcolor{red}{разомкнутого} треугольника»}.  Данный  переключатель  позволяет  учесть  различные  схемы  соединения дополнительной  обмотки  ТН.  Если  переключатель  выставлен  в  положение  «$U_{\textsl{НК}}$»,  то  переменной  $3U_{\textsl{0 треуг.}}$ присваивается  значение  напряжения,  измеренного  по  аналоговому  входу  «$U_{\textsl{ИФ}}(U_{\textsl{ИК}})$».  Если  программный переключатель \textcolor{red}{\textbf{«Схема подключения цепей треугольника»}} выставлен в положение «$U_{\textsl{НИ}}/U_{\textsl{ИК}}$», «$U_{\textsl{НФ}}/U_{\textsl{ФК}}$» или  «$U_{\textsl{НИ}}/U_{\textsl{ИФ}}/U_{\textsl{ФК}}$», то переменная $3U_{\textsl{0 треуг.}}$ рассчитывается в соответствии с формулами, приведенными в \textcolor{red}{приложении А}. Программный переключатель \textbf{«Выбор 3U0»} аналогично влияет на работу ОНМНП. 

\item
Устройство содержит шесть ступеней ДЗ от междуфазных замыканий с контролем трех контуров «фаза-фаза», а также три ступени с контролем контуров «фаза-земля» для защиты от однофазных КЗ. 

\end{enumerate}

%%%%%%%%%%%%%%%%%%%%%%%%%%%%%%%%%%%%%%%%%%%%%%%%%%%%%%%%%%%%%%%%%%%%%% ПЕРВАЯ СТУПЕНЬ %%%%%%%%%%%%%%%%%%%%%%%%%%%%%%%%%%%%%%%%%%%%%%%%%%
\item\label{sec:sec2} Первая ступень ДЗ от междуфазных КЗ (ДЗ-ФФ-1)
\begin{enumerate}[label=\arabic{section}.\arabic{subsection}.\arabic{enumi}.\arabic*, labelsep=4pt, leftmargin=0em, itemindent=65pt, parsep=0pt]

\item
Первая  ступень  ДЗ  от  междуфазных  КЗ  (ДЗ-ФФ-1)  содержит  две  «подступени»:  быстродействующая ДЗ-ФФ-1(б)  и  медленнодействующая  ДЗ-ФФ-1(м),  которая  вводится программным  ключом  «ДЗ-ФФ-1(м)». Всю ступень можно ввести/вывести программным ключом \textbf{«ДЗ-ФЗ-1»}. Контроль от блокировки при качаниях (БК)  для  каждой  подступени  задается  отдельными  программными  переключателями  \textbf{«БК  ДЗ-ФФ-1(б)»}  и  \textbf{«БК  ДЗ-ФФ-1(м)»}.  Действие  подступеней  осуществляется  с  регулируемыми  независимыми  выдержками времени \textbf{«Тср ДЗ-ФФ-1(б)»} и \textbf{«Тср ДЗ-ФФ-1(м)»}.

\item
Предусмотрена возможность подхвата сработанного состояния РС-ФФ-1 от направленного РС-ФФ-2 с охватом начала координат (вводится программным ключом \textbf{«Подхват РС-ФФ-1»}). Тем самым гарантируется правильное действие защиты при близких  КЗ.  Возврат схемы подхвата в исходное состояние происходит только после возврата направленного РС-ФФ-2 с охватом начала координат. 

\item
Если междуфазное КЗ произошло в зоне действия ДЗ-ФФ-1 и при этом в течение времени ввода БК сработало направленное РС-ФФ-2 с охватом начала координат, то разрешающий сигнал от БК подхватывается и удерживается независимо от истечения времени ввода. Возврат разрешающего сигнала происходит только после возврата направленного РС-ФФ-2 с охватом начала координат. 

\item
Функциональная схема первой ступени  ДЗ  от  междуфазных  КЗ приведена на рисунке \ref{pict_dz:ff1}.

\begin{figure}[H]
\centering
\includegraphics[width=1\textwidth,height=1\textheight,keepaspectratio]{\fbpath/561.0021 ДЗ/img/_ff1.pdf}
\captionof{figure}{Функциональная схема ДЗ-ФФ-1}\label{pict_dz:ff1}
\end{figure}

\end{enumerate}
%%%%%%%%%%%%%%%%%%%%%%%%%%%%%%%%%%%%%%%%%%%%%%%%%%%%%%%%%%%%%%%%%%%%%% КОНЕЦ ПЕРВАЯ СТУПЕНЬ %%%%%%%%%%%%%%%%%%%%%%%%%%%%%%%%%%%%%%%%%%%%%%%%%%

%%%%%%%%%%%%%%%%%%%%%%%%%%%%%%%%%%%%%%%%%%%%%%%%%%%%%%%%%%%%%%%%%%%%%% ВТОРАЯ СТУПЕНЬ %%%%%%%%%%%%%%%%%%%%%%%%%%%%%%%%%%%%%%%%%%%%%%%%%%
\item\label{sec:sec2} Вторая ступень ДЗ от междуфазных КЗ (ДЗ-ФФ-1)
\begin{enumerate}[label=\arabic{section}.\arabic{subsection}.\arabic{enumi}.\arabic*, labelsep=4pt, leftmargin=0em, itemindent=65pt, parsep=0pt]

\item
Вторая ступень ДЗ от междуфазных КЗ также выполнена с двумя выдержками времени и с раздельным контролем от БК. Вся ступень вводится программным ключом \textbf{«ДЗ-ФФ-2»}. Быстродействующую подступень можно  ввести  в  работу  программным  ключом  \textbf{«ДЗ-ФФ-2(б)»}.  Действие  подступеней  осуществляется  с регулируемыми независимыми выдержками времени \textbf{«Тср ДЗ-ФФ-2(б)»} и \textbf{«Тср ДЗ-ФФ-2(м)»}. Контроль от БК для каждой подступени задается отдельными программными переключателями \textbf{«БК ДЗ-ФФ-2(б)»} и \textbf{«БК ДЗ-ФФ-2(м)»}.

\item
Разрешающий сигнал от БК для ДЗ-ФФ-2 подхватывается от направленного РС-ФФ-2 с охватом начала координат  и  удерживается  независимо  от  истечения  времени  ввода  БК.  Возврат  разрешающего  сигнала происходит только после возврата направленного РС-ФФ-2 с охватом начала координат. 

\begin{figure}[H]
\centering
\includegraphics[width=1\textwidth,height=1\textheight,keepaspectratio]{\fbpath/561.0021 ДЗ/img/_ff2.pdf}
\captionof{figure}{Функциональная схема ДЗ-ФФ-2}\label{pict_dz:ff2}
\end{figure}

\end{enumerate}
%%%%%%%%%%%%%%%%%%%%%%%%%%%%%%%%%%%%%%%%%%%%%%%%%%%%%%%%%%%%%%%%%%%%%% КОНЕЦ ВТОРАЯ СТУПЕНЬ %%%%%%%%%%%%%%%%%%%%%%%%%%%%%%%%%%%%%%%%%%%%%%%%%%


%%%%%%%%%%%%%%%%%%%%%%%%%%%%%%%%%%%%%%%%%%%%%%%%%%%%%%%%%%%%%%%%%%%%%% 3,4,5,6 СТУПЕНИ %%%%%%%%%%%%%%%%%%%%%%%%%%%%%%%%%%%%%%%%%%%%%%%%%%
\item\label{sec:sec2} Остальные ступени ДЗ от междуфазных КЗ (ДЗ-ФФ-3, ДЗ-ФФ-4, ДЗ-ФФ-5, ДЗ-ФФ-6)
\begin{enumerate}[label=\arabic{section}.\arabic{subsection}.\arabic{enumi}.\arabic*, labelsep=4pt, leftmargin=0em, itemindent=65pt, parsep=0pt]


\item
Логика  работы  остальных  ступеней  ДЗ  от  междуфазных  КЗ  выполнена  одинаковой.  Ввод  в  работу ступеней осуществляется программными ключами \textbf{«ДЗ-ФФ-3(4,5,6)»}. Действие ступеней осуществляется с регулируемыми  независимыми  выдержками  времени  \textbf{«Тср  ДЗ-ФФ-3(4,5,6)»}.  Контроль  от  БК  задается отдельными программными переключателями \textbf{«БК ДЗ-ФФ-3(4,5,6)»}.  

\item
Разрешающий  сигнал  от  БК  для  ступеней  подхватывается  от  срабатывания  РС  соответствующей ступени  и  удерживается  независимо  от  истечения  времени  ввода  БК.  Возврат  разрешающего  сигнала происходит только после возврата РС соответствующей ступени. На рисунке 9 приведена функциональная схема ДЗ-ФФ-3, выполнение функциональных схем ступеней ДЗ ДЗ-ФФ-4, ДЗ-ФФ-5, ДЗ-ФФ-6 аналогичное. 
























\item
Уставки  <<Вырез  нагрузки>>,  <<Rнг>>,  <<Фнг>>  являются  общими  для  всех  ступеней ДЗ, а уставка SGF10 общая только для ступеней ДЗ <<фаза-земля>>. Общие параметры для настройки функции <<ДЗ>> приведены в таблице \ref{dz:tbl0}.

%================================================== Таблица =====================================
\small
\begin{longtable}{|>{\centering\arraybackslash}m{5.3cm}|>{\centering\arraybackslash}m{3.3cm}|>{\centering\arraybackslash}m{4.2cm}|>{\centering\arraybackslash}m{1.8cm}|>{\centering\arraybackslash}m{1cm}|}
\caption{Общие параметры для настройки функции <<ДЗ>>\hfill\vspace{-0.5\baselineskip}}\label{dz:tbl0}\\ 
\hline
\rowcolor{gray!30}
Параметр (Параметр на ИЧМ) & Условное обозначение на схеме & Значение/ Диапазон & Единица измерения & Шаг \\ 
\hline
\endfirsthead
\caption*{\hspace{3pt}\emph{Продолжение таблицы \ref{dz:tbl0}\hfill\vspace{-0.5\baselineskip}}} \\ % сделано по ГОСТ 2.105 п.6.8.7
\hline
\rowcolor{gray!30}
Параметр (Параметр на ИЧМ) & Условное обозначение на схеме & Значение/ Диапазон & Единица измерения & Шаг \\ 
%\hline
\endhead
%\hline
\endfoot
%\hline
\endlastfoot
%==+t1*LLN0|LVDIS
\centering Учет нагрузочного режима в характеристике срабатывания (Вырез\_нагрузки) & \centering SGF2 & \centering 0 = Не предусмотрено\\1 = Предусмотрено & \centering -- & \centering \arraybackslash -- \\
\hline
\centering Активное сопротивление для отстройки от режима максимальной нагрузки (Rнг) & \centering -- & \centering 0,01 ... 500,00 & \centering Ом & \centering \arraybackslash 0,01 \\
\hline
\centering Угол наклона для отстройки от режима максимальной нагрузки (Фнг) & \centering -- & \centering 5,0 ... 60,0 & \centering град. & \centering \arraybackslash 0,1 \\
\hline
\centering Режим работы ДЗ по контуру «фаза В - земля» (Ввод\_B0) & \centering SGF10 & \centering 0 = Не предусмотрено\\1 = Предусмотрено & \centering -- & \centering \arraybackslash -- \\
\hline
\centering Удельное активное сопротивление прямой последовательности (R1уд) & \centering -- & \centering 0,001 ... 100,000 & \centering Ом/км & \centering \arraybackslash 0,001 \\
\hline
\centering Удельное реактивное сопротивление прямой последовательности (X1уд) & \centering -- & \centering 0,001 ... 100,000 & \centering Ом/км & \centering \arraybackslash 0,001 \\
\hline
\centering Удельное активное сопротивление нулевой последовательности (R0уд) & \centering -- & \centering 0,001 ... 100,000 & \centering Ом/км & \centering \arraybackslash 0,001 \\
\hline
\centering Удельное реактивное сопротивление нулевой последовательности (X0уд) & \centering -- & \centering 0,001 ... 100,000 & \centering Ом/км & \centering \arraybackslash 0,001 \\
\hline
%===t1
\end{longtable}
\normalsize
%================================================== Конец таблицы =====================================

\item
Ввод каждой ступени в работу на этапе параметрирования устройства осуществляется программным переключателем SGF1 <<Ввод\_функции>>. Информация о введенной ступени отображается выходным сигналом <<ДЗ-фф(фз) (n)ст.: Ввод>>.
\item
Каждая ступень <<ДЗ>> может быть выведена из работы оперативно путем активации сигнала <<ОВ ДЗ-фф(фз) (n)ст.>>, что характеризуется активным состоянием сигнала <<ДЗ-фф(фз) (n)ст.: Оперативный вывод>> на выходе функции.

\item
При активном состоянии входного сигнала <<ДЗ-фф(фз)(n)ст. на сигн>> действие ступени переводится на сигнал.

\item
Алгоритм работы первой ступени ДЗ по контуру <<фаза-фаза>> (<<ДЗ-фф 1ст.>>) выполнен в соответствии с рисунком \ref{dz:dz1ffimg1}. В таблице \ref{dz:dz1fftbl1} приведены параметры, необходимые для настройки функции <<ДЗ-фф 1ст.>>.

\vspace{3mm}
\begin{figure}[H]
\centering
\includegraphics[width=0.9\textwidth,height=0.9\textheight,keepaspectratio]{\fbpath/561.0021 ДЗ/img/dz1ff.pdf}
\captionof{figure}{Функциональная схема <<ДЗ-фф 1ст.>>}\label{dz:dz1ffimg1}
\end{figure}

\small
\begin{longtable}{|>{\centering\arraybackslash}m{5.3cm}|>{\centering\arraybackslash}m{3.3cm}|>{\centering\arraybackslash}m{4.2cm}|>{\centering\arraybackslash}m{1.8cm}|>{\centering\arraybackslash}m{1cm}|}
\caption{Параметры для настройки функции <<ДЗ-фф 1ст.>>\hfill\vspace{-0.5\baselineskip}}\label{dz:dz1fftbl1}\\ 
\hline
\rowcolor{gray!30}
Параметр (Параметр на ИЧМ) & Условное обозначение на схеме & Значение/ Диапазон & Единица измерения & Шаг \\ 
\hline
\endfirsthead
\caption*{\hspace{3pt}\emph{Продолжение таблицы \ref{dz:dz1fftbl1}\hfill\vspace{-0.5\baselineskip}}} \\ % сделано по ГОСТ 2.105 п.6.8.7
\hline
\rowcolor{gray!30}
Параметр (Параметр на ИЧМ) & Условное обозначение на схеме & Значение/ Диапазон & Единица измерения & Шаг \\ 
%\hline
\endhead
%\hline
\endfoot
%\hline
\endlastfoot
%==+t1*PPPDIS1|LVDIS
\centering Ввод функции в работу (Ввод\_функции) & \centering SGF1 & \centering 0 = Не предусмотрено\\1 = Предусмотрено & \centering -- & \centering \arraybackslash -- \\
\hline
\centering Угол наклона верхней стороны характеристики срабатывания для компенсации нагрузки (Фкн) & \centering -- & \centering -45,0 ... 0,0 & \centering град. & \centering \arraybackslash 0,1 \\
\hline
\centering Активное сопротивление срабатывания (Rср) & \centering -- & \centering 0,01 ... 500,00 & \centering Ом & \centering \arraybackslash 0,01 \\
\hline
\centering Реактивное сопротивление срабатывания (Xср) & \centering -- & \centering 0,01 ... 500,00 & \centering Ом & \centering \arraybackslash 0,01 \\
\hline
\centering Угол наклона характеристики срабатывания (Фхс) & \centering -- & \centering 30,0 ... 89,9 & \centering град. & \centering \arraybackslash 0,1 \\
\hline
\centering Режим направленности (Направленность) & \centering SGF3 & \centering 1 = Ненаправленная\\2 = Прямонаправленная\\3 = Обратнонаправленная & \centering -- & \centering \arraybackslash -- \\
\hline
\centering Подхват от ненаправленного пуска или отдельной ступени  (Подхват) & \centering SGF9 & \centering 0 = Не предусмотрено\\1 = Предусмотрено & \centering -- & \centering \arraybackslash -- \\
\hline
\centering Режим контроля от БК (Реж\_БК) & \centering SGF7 & \centering 1 = С контролем от БК-б\\2 = С контролем от БК-м & \centering -- & \centering \arraybackslash -- \\
\hline
\centering Режим контроля от БНН (Реж\_БНН) & \centering SGF8 & \centering 0 = Без контроля\\1 = С контролем & \centering -- & \centering \arraybackslash -- \\
\hline
\centering Режим контроля от ОПО (Реж\_ОПО) & \centering SGF6 & \centering 1 = По току\\2 = По току и напряжению\\3 = По току или по току и напряжению & \centering -- & \centering \arraybackslash -- \\
\hline
\centering Тип ПО (Контр\_ПО) & \centering SGF5 & \centering 1 = Без контроля\\2 = С контролем от БК\\3 = С контролем от ОПО & \centering -- & \centering \arraybackslash -- \\
\hline
\centering Вывод направленности при АУ (Ненапр\_АУ) & \centering SGF4 & \centering 0 = Не предусмотрено\\1 = Предусмотрено & \centering -- & \centering \arraybackslash -- \\
\hline
\centering Выдержка времени срабатывания (Tср) & \centering T1 & \centering 0,00 ... 10,00 & \centering с & \centering \arraybackslash 0,01 \\
\hline
%===t1
\end{longtable}
\normalsize

\item
Алгоритм работы второй ступени ДЗ по контуру <<фаза-фаза>> (<<ДЗ-фф 2ст.>>) выполнен в соответствии с рисунком \ref{dz:dz2ffimg1}. В таблице \ref{dz:dz2fftbl1} приведены параметры, необходимые для настройки функции <<ДЗ-фф 2ст.>>.

\small
\begin{longtable}{|>{\centering\arraybackslash}m{5.3cm}|>{\centering\arraybackslash}m{3.3cm}|>{\centering\arraybackslash}m{4.2cm}|>{\centering\arraybackslash}m{1.8cm}|>{\centering\arraybackslash}m{1cm}|}
\caption{Параметры для настройки функции <<ДЗ-фф 2ст.>>\hfill\vspace{-0.5\baselineskip}}\label{dz:dz2fftbl1}\\ 
\hline
\rowcolor{gray!30}
Параметр (Параметр на ИЧМ) & Условное обозначение на схеме & Значение/ Диапазон & Единица измерения & Шаг \\ 
\hline
\endfirsthead
%\caption{\emph{Продолжение\hfill\vspace{-0.5\baselineskip}}} \\ % Отступ обозначения таблицы от таблицы и выравнивание надписи слева
\caption*{\hspace{3pt}\emph{Продолжение таблицы \ref{dz:dz2fftbl1}\hfill\vspace{-0.5\baselineskip}}} \\ % сделано по ГОСТ 2.105 п.6.8.7
\hline
\rowcolor{gray!30}
Параметр (Параметр на ИЧМ) & Условное обозначение на схеме & Значение/ Диапазон & Единица измерения & Шаг \\ 
%\hline
\endhead
%\hline
\endfoot
%\hline
\endlastfoot
%==+t1*PPPDIS2|LVDIS
\centering Ввод функции в работу (Ввод\_функции) & \centering SGF1 & \centering 0 = Не предусмотрено\\1 = Предусмотрено & \centering -- & \centering \arraybackslash -- \\
\hline
\centering Угол наклона верхней стороны характеристики срабатывания для компенсации нагрузки (Фкн) & \centering -- & \centering -45,0 ... 0,0 & \centering град. & \centering \arraybackslash 0,1 \\
\hline
\centering Активное сопротивление срабатывания (Rср) & \centering -- & \centering 0,01 ... 500,00 & \centering Ом & \centering \arraybackslash 0,01 \\
\hline
\centering Реактивное сопротивление срабатывания (Xср) & \centering -- & \centering 0,01 ... 500,00 & \centering Ом & \centering \arraybackslash 0,01 \\
\hline
\centering Угол наклона характеристики срабатывания (Фхс) & \centering -- & \centering 30,0 ... 89,9 & \centering град. & \centering \arraybackslash 0,1 \\
\hline
\centering Режим направленности (Направленность) & \centering SGF3 & \centering 1 = Ненаправленная\\2 = Прямонаправленная\\3 = Обратнонаправленная & \centering -- & \centering \arraybackslash -- \\
\hline
\centering Режим контроля от БК (Реж\_БК) & \centering SGF7 & \centering 1 = С контролем от БК-б\\2 = С контролем от БК-м & \centering -- & \centering \arraybackslash -- \\
\hline
\centering Режим контроля от БНН (Реж\_БНН) & \centering SGF8 & \centering 0 = Без контроля\\1 = С контролем & \centering -- & \centering \arraybackslash -- \\
\hline
\centering Режим контроля от ОПО (Реж\_ОПО) & \centering SGF6 & \centering 1 = По току\\2 = По току и напряжению\\3 = По току или по току и напряжению & \centering -- & \centering \arraybackslash -- \\
\hline
\centering Тип ПО (Контр\_ПО) & \centering SGF5 & \centering 1 = Без контроля\\2 = С контролем от БК\\3 = С контролем от ОПО & \centering -- & \centering \arraybackslash -- \\
\hline
\centering Вывод направленности при АУ (Ненапр\_АУ) & \centering SGF4 & \centering 0 = Не предусмотрено\\1 = Предусмотрено & \centering -- & \centering \arraybackslash -- \\
\hline
\centering Выдержка времени срабатывания (Tср) & \centering T1 & \centering 0,00 ... 10,00 & \centering с & \centering \arraybackslash 0,01 \\
\hline
%===t1
\end{longtable}
\normalsize

\vspace{3mm}
\begin{figure}[H]
\centering
\includegraphics[width=0.9\textwidth,height=0.9\textheight,keepaspectratio]{\fbpath/561.0021 ДЗ/img/dz2ff.pdf}
\captionof{figure}{Функциональная схема <<ДЗ-фф 2ст.>>}\label{dz:dz2ffimg1}
\end{figure}

\item
Алгоритм работы третьей ступени ДЗ по контуру <<фаза-фаза>> (<<ДЗ-фф 3ст.>>) выполнен в соответствии с рисунком \ref{dz:dz3ffimg1}. В таблице \ref{dz:dz3fftbl1} приведены параметры, необходимые для настройки функции <<ДЗ-фф 3ст.>>.

\small
\begin{longtable}{|>{\centering\arraybackslash}m{5.3cm}|>{\centering\arraybackslash}m{3.3cm}|>{\centering\arraybackslash}m{4.2cm}|>{\centering\arraybackslash}m{1.8cm}|>{\centering\arraybackslash}m{1cm}|}
\caption{Параметры для настройки функции <<ДЗ-фф 3ст.>>\hfill\vspace{-0.5\baselineskip}}\label{dz:dz3fftbl1}\\ 
\hline
\rowcolor{gray!30}
Параметр (Параметр на ИЧМ) & Условное обозначение на схеме & Значение/ Диапазон & Единица измерения & Шаг \\ 
\hline
\endfirsthead
\caption*{\hspace{3pt}\emph{Продолжение таблицы \ref{dz:dz3fftbl1}\hfill\vspace{-0.5\baselineskip}}} \\ % сделано по ГОСТ 2.105 п.6.8.7
\hline
\rowcolor{gray!30}
Параметр (Параметр на ИЧМ) & Условное обозначение на схеме & Значение/ Диапазон & Единица измерения & Шаг \\ 
%\hline
\endhead
%\hline
\endfoot
%\hline
\endlastfoot
%==+t1*PPPDIS3|LVDIS
\centering Ввод функции в работу (Ввод\_функции) & \centering SGF1 & \centering 0 = Не предусмотрено\\1 = Предусмотрено & \centering -- & \centering \arraybackslash -- \\
\hline
\centering Угол наклона верхней стороны характеристики срабатывания для компенсации нагрузки (Фкн) & \centering -- & \centering -45,0 ... 0,0 & \centering град. & \centering \arraybackslash 0,1 \\
\hline
\centering Активное сопротивление срабатывания (Rср) & \centering -- & \centering 0,01 ... 500,00 & \centering Ом & \centering \arraybackslash 0,01 \\
\hline
\centering Реактивное сопротивление срабатывания (Xср) & \centering -- & \centering 0,01 ... 500,00 & \centering Ом & \centering \arraybackslash 0,01 \\
\hline
\centering Угол наклона характеристики срабатывания (Фхс) & \centering -- & \centering 30,0 ... 89,9 & \centering град. & \centering \arraybackslash 0,1 \\
\hline
\centering Режим направленности (Направленность) & \centering SGF3 & \centering 1 = Ненаправленная\\2 = Прямонаправленная\\3 = Обратнонаправленная & \centering -- & \centering \arraybackslash -- \\
\hline
\centering Режим контроля от БК (Реж\_БК) & \centering SGF7 & \centering 1 = С контролем от БК-б\\2 = С контролем от БК-м & \centering -- & \centering \arraybackslash -- \\
\hline
\centering Режим контроля от БНН (Реж\_БНН) & \centering SGF8 & \centering 0 = Без контроля\\1 = С контролем & \centering -- & \centering \arraybackslash -- \\
\hline
\centering Режим контроля от ОПО (Реж\_ОПО) & \centering SGF6 & \centering 1 = По току\\2 = По току и напряжению\\3 = По току или по току и напряжению & \centering -- & \centering \arraybackslash -- \\
\hline
\centering Тип ПО (Контр\_ПО) & \centering SGF5 & \centering 1 = Без контроля\\2 = С контролем от БК\\3 = С контролем от ОПО & \centering -- & \centering \arraybackslash -- \\
\hline
\centering Вывод направленности при АУ (Ненапр\_АУ) & \centering SGF4 & \centering 0 = Не предусмотрено\\1 = Предусмотрено & \centering -- & \centering \arraybackslash -- \\
\hline
\centering Выдержка времени срабатывания (Tср) & \centering T1 & \centering 0,00 ... 10,00 & \centering с & \centering \arraybackslash 0,01 \\
\hline
%===t1
\end{longtable}
\normalsize

\vspace{3mm}
\begin{figure}[H]
\centering
\includegraphics[width=0.9\textwidth,height=0.9\textheight,keepaspectratio]{\fbpath/561.0021 ДЗ/img/dz3ff.pdf}
\captionof{figure}{Функциональная схема <<ДЗ-фф 3ст.>>}\label{dz:dz3ffimg1}
\end{figure}

\item
Алгоритм работы первой ступени ДЗ по контуру <<фаза-земля>> (<<ДЗ-фз 1ст.>>) выполнен в соответствии с рисунком \ref{dz:dz1fzimg1}. В таблице \ref{dz:dz1fztbl1} приведены параметры, необходимые для настройки функции <<ДЗ-фз 1ст.>>.

\small
\begin{longtable}{|>{\centering\arraybackslash}m{5.3cm}|>{\centering\arraybackslash}m{3.3cm}|>{\centering\arraybackslash}m{4.2cm}|>{\centering\arraybackslash}m{1.8cm}|>{\centering\arraybackslash}m{1cm}|}
\caption{Параметры для настройки функции <<ДЗ-фз 1ст.>>\hfill\vspace{-0.5\baselineskip}}\label{dz:dz1fztbl1}\\ 
\hline
\rowcolor{gray!30}
Параметр (Параметр на ИЧМ) & Условное обозначение на схеме & Значение/ Диапазон & Единица измерения & Шаг \\ 
\hline
\endfirsthead
\caption*{\hspace{3pt}\emph{Продолжение таблицы \ref{dz:dz1fztbl1}\hfill\vspace{-0.5\baselineskip}}} \\ % сделано по ГОСТ 2.105 п.6.8.7
\hline
\rowcolor{gray!30}
Параметр (Параметр на ИЧМ) & Условное обозначение на схеме & Значение/ Диапазон & Единица измерения & Шаг \\ 
%\hline
\endhead
%\hline
\endfoot
%\hline
\endlastfoot
%==+t1*PEPDIS1|LVDIS
\centering Ввод функции в работу (Ввод\_функции) & \centering SGF1 & \centering 0 = Не предусмотрено\\1 = Предусмотрено & \centering -- & \centering \arraybackslash -- \\
\hline
\centering Угол наклона верхней стороны характеристики срабатывания для компенсации нагрузки (Фкн) & \centering -- & \centering -45,0 ... 0,0 & \centering град. & \centering \arraybackslash 0,1 \\
\hline
\centering Активное сопротивление срабатывания (Rср) & \centering -- & \centering 0,01 ... 500,00 & \centering Ом & \centering \arraybackslash 0,01 \\
\hline
\centering Реактивное сопротивление срабатывания (Xср) & \centering -- & \centering 0,01 ... 500,00 & \centering Ом & \centering \arraybackslash 0,01 \\
\hline
\centering Угол наклона характеристики срабатывания (Фхс) & \centering -- & \centering 30,0 ... 89,9 & \centering град. & \centering \arraybackslash 0,1 \\
\hline
\centering Режим направленности (Направленность) & \centering SGF3 & \centering 1 = Ненаправленная\\2 = Прямонаправленная\\3 = Обратнонаправленная & \centering -- & \centering \arraybackslash -- \\
\hline
\centering Подхват от ненаправленного пуска или отдельной ступени  (Подхват) & \centering SGF9 & \centering 0 = Не предусмотрено\\1 = Предусмотрено & \centering -- & \centering \arraybackslash -- \\
\hline
\centering Режим контроля от БК (Реж\_БК) & \centering SGF7 & \centering 1 = С контролем от БК-б\\2 = С контролем от БК-м & \centering -- & \centering \arraybackslash -- \\
\hline
\centering Режим контроля от БНН (Реж\_БНН) & \centering SGF8 & \centering 0 = Без контроля\\1 = С контролем & \centering -- & \centering \arraybackslash -- \\
\hline
\centering Режим контроля от ОПО (Реж\_ОПО) & \centering SGF6 & \centering 1 = По току\\2 = По току и напряжению\\3 = По току или по току и напряжению & \centering -- & \centering \arraybackslash -- \\
\hline
\centering Тип ПО (Контр\_ПО) & \centering SGF5 & \centering 1 = Без контроля\\2 = С контролем от БК\\3 = С контролем от ОПО & \centering -- & \centering \arraybackslash -- \\
\hline
\centering Вывод направленности при АУ (Ненапр\_АУ) & \centering SGF4 & \centering 0 = Не предусмотрено\\1 = Предусмотрено & \centering -- & \centering \arraybackslash -- \\
\hline
\centering Выдержка времени срабатывания (Tср) & \centering T1 & \centering 0,00 ... 10,00 & \centering с & \centering \arraybackslash 0,01 \\
\hline
%===t1
\end{longtable}
\normalsize

\vspace{3mm}
\begin{figure}[H]
\centering
\includegraphics[width=0.9\textwidth,height=0.9\textheight,keepaspectratio]{\fbpath/561.0021 ДЗ/img/dz1fz.pdf}
\captionof{figure}{Функциональная схема <<ДЗ-фз 1ст.>>}\label{dz:dz1fzimg1}
\end{figure}

\item
Алгоритм работы второй ступени ДЗ по контуру <<фаза-земля>> (<<ДЗ-фз 2ст.>>) выполнен в соответствии с рисунком \ref{dz:dz2fzimg1}. В таблице \ref{dz:dz2fztbl1} приведены параметры, необходимые для настройки функции <<ДЗ-фз 2ст.>>.

\small
\begin{longtable}{|>{\centering\arraybackslash}m{5.3cm}|>{\centering\arraybackslash}m{3.3cm}|>{\centering\arraybackslash}m{4.2cm}|>{\centering\arraybackslash}m{1.8cm}|>{\centering\arraybackslash}m{1cm}|}
\caption{Параметры для настройки функции <<ДЗ-фз 2ст.>>\hfill\vspace{-0.5\baselineskip}}\label{dz:dz2fztbl1}\\ 
\hline
\rowcolor{gray!30}
Параметр (Параметр на ИЧМ) & Условное обозначение на схеме & Значение/ Диапазон & Единица измерения & Шаг \\ 
\hline
\endfirsthead
\caption*{\hspace{3pt}\emph{Продолжение таблицы \ref{dz:dz2fztbl1}\hfill\vspace{-0.5\baselineskip}}} \\ % сделано по ГОСТ 2.105 п.6.8.7
\hline
\rowcolor{gray!30}
Параметр (Параметр на ИЧМ) & Условное обозначение на схеме & Значение/ Диапазон & Единица измерения & Шаг \\ 
%\hline
\endhead
%\hline
\endfoot
%\hline
\endlastfoot
%==+t1*PEPDIS2|LVDIS
\centering Ввод функции в работу (Ввод\_функции) & \centering SGF1 & \centering 0 = Не предусмотрено\\1 = Предусмотрено & \centering -- & \centering \arraybackslash -- \\
\hline
\centering Угол наклона верхней стороны характеристики срабатывания для компенсации нагрузки (Фкн) & \centering -- & \centering -45,0 ... 0,0 & \centering град. & \centering \arraybackslash 0,1 \\
\hline
\centering Активное сопротивление срабатывания (Rср) & \centering -- & \centering 0,01 ... 500,00 & \centering Ом & \centering \arraybackslash 0,01 \\
\hline
\centering Реактивное сопротивление срабатывания (Xср) & \centering -- & \centering 0,01 ... 500,00 & \centering Ом & \centering \arraybackslash 0,01 \\
\hline
\centering Угол наклона характеристики срабатывания (Фхс) & \centering -- & \centering 30,0 ... 89,9 & \centering град. & \centering \arraybackslash 0,1 \\
\hline
\centering Режим направленности (Направленность) & \centering SGF3 & \centering 1 = Ненаправленная\\2 = Прямонаправленная\\3 = Обратнонаправленная & \centering -- & \centering \arraybackslash -- \\
\hline
\centering Режим контроля от БК (Реж\_БК) & \centering SGF7 & \centering 1 = С контролем от БК-б\\2 = С контролем от БК-м & \centering -- & \centering \arraybackslash -- \\
\hline
\centering Режим контроля от БНН (Реж\_БНН) & \centering SGF8 & \centering 0 = Без контроля\\1 = С контролем & \centering -- & \centering \arraybackslash -- \\
\hline
\centering Режим контроля от ОПО (Реж\_ОПО) & \centering SGF6 & \centering 1 = По току\\2 = По току и напряжению\\3 = По току или по току и напряжению & \centering -- & \centering \arraybackslash -- \\
\hline
\centering Тип ПО (Контр\_ПО) & \centering SGF5 & \centering 1 = Без контроля\\2 = С контролем от БК\\3 = С контролем от ОПО & \centering -- & \centering \arraybackslash -- \\
\hline
\centering Вывод направленности при АУ (Ненапр\_АУ) & \centering SGF4 & \centering 0 = Не предусмотрено\\1 = Предусмотрено & \centering -- & \centering \arraybackslash -- \\
\hline
\centering Выдержка времени срабатывания (Tср) & \centering T1 & \centering 0,00 ... 10,00 & \centering с & \centering \arraybackslash 0,01 \\
\hline
%===t1
\end{longtable}
\normalsize

\vspace{3mm}
\begin{figure}[H]
\centering
\includegraphics[width=0.9\textwidth,height=0.9\textheight,keepaspectratio]{\fbpath/561.0021 ДЗ/img/dz2fz.pdf}
\captionof{figure}{Функциональная схема <<ДЗ-фз 2ст.>>}\label{dz:dz2fzimg1}
\end{figure}

\end{enumerate}

%%%%%%%%%%%%%%%%%%%%%%%%%%%%%%%%%%%%%%%%%%%%%%%%%%%%%%%%%%%%%%%%%%%%%% ОРГАН НАПРАВЛЕНИЯ %%%%%%%%%%%%%%%%%%%%%%%%%%%%%%%%%%%%%%%%%%%%%%%%%%
\item\label{sec:on} Орган направления (ОН)
\begin{enumerate}[label=\arabic{section}.\arabic{subsection}.\arabic{enumi}.\arabic*, labelsep=4pt, leftmargin=0em, itemindent=65pt, parsep=0pt]

\item
В кольцевых сетях, в сетях с двухсторонним питанием и особенно в сложных сетях с несколькими источниками питания в большинстве случаев ненаправленная защита не может обеспечить селективность действия. Для  селективной  работы в таких случаях необходимо применение органа  направления мощности совместно с измерительным токовым органом защиты. Защита действует на отключение только в  том  случае, если срабатывает как токовый измерительный орган, так и орган  направления, который разрешает действие при направлении мощности КЗ от шин в линию.
\item
Функция  <<ОН>>  становится  активна,  если  введена  в  работу  одна  из  ступеней  ДЗ и в ней настроен контроль направленности. 
\item
ОН  применим  для  реализации  направленных  ступеней  ДЗ.  Направленная  ступень  ДЗ  будет действовать  только  в  том  случае,  когда  сработает  реле  сопротивления,  а  также  ОН,  разрешающий работать ДЗ в прямом направлении. Функция ОН содержит в себе шесть органов, определяющих направление (три для фазных органов ДЗ, три для междуфазных органов ДЗ). 
\item
ОН  выполнен  в  виде  отдельного  функционального  модуля  с  собственными  уставками.  ОН  ДЗ осуществляет построение характеристики ОН по уставкам, ограничивающим характеристику срабатывания органа в прямом направлении во втором и четвертом квадрантах (<<Угол во 2 квадр.>> - $\varphi2$ и <<Угол в 4 квадр.>> – $\varphi1$ соответственно), и, используя положение  вектора  сопротивления  на  плоскости  характеристики, формирует выходные сигналы прямого либо обратного направления согласно характеристике, изображенной на \ref{dz:onimg0}.

\vspace{3mm}
\begin{figure}[H]
\centering
\includegraphics[width=0.5\textwidth,height=0.5\textheight,keepaspectratio]{\fbpath/561.0021 ДЗ/img/charon.pdf}
\captionof{figure}{{Характеристика ОН}}\label{dz:onimg0}
\end{figure}

\item
ОН ДЗ может работать в одном из режимов направленности:
\begin{itemize} 
    \item в прямом направлении, если вектор сопротивления находится в диапазоне от $\varphi1$ до $\varphi2$; 
    \item в обратном направлении, если вектор сопротивления находится в диапазоне от $\varphi1-180^\circ$ до $\varphi2-180^\circ$; 
\item без направления, если вектор сопротивления находится в области нечувствительности.
\end{itemize}

\item
Строгая направленность обеспечивается действиями ступеней ДЗ при КЗ вблизи установки защиты используется специальный ОН на аварийных составляющих, который вводится при снижении линейного напряжения ниже 5 В.
\item
Данный орган контролирует направление повреждения и блокирует срабатывание ступеней ДЗ при КЗ <<за спиной>>. ОН на аварийных составляющих выполнен таким образом, что обеспечивает правильное 
определение  направления  повреждения  при  любых  видах  КЗ  (в  том  числе  при  близких  с  просадкой напряжения).  Обеспечивается  отсутствие  ложных  срабатываний  ОН  при  КЗ  <<за  спиной>>  при  токах  до $20 \cdot I_{\textsl{ном}}$.
\item
Отказ ОН возможен при включении на КЗ, если защита подключена к ТН, установленного на линии. Эта  особенность  характерна  всем  известным  способам,  работа  которых  опирается  на  напряжение 
предшествующего режима.
\item
Принцип работы фазных органов следующий:
\begin{itemize} 
    \item в качестве воздействующих величин используются действующие значения первой гармоники фазных токов и значения расчетных напряжений поляризации, рассчитанные по формулам (\ref{eq:group1}); 
    \item разрешающими  условиями  для  работы  направленного  органа  являются  превышение  тока значения  $0,04 \cdot I_{\textsl{ном}}$  и превышение напряжения значения 1 вольт; 
\item фаза вектора сопротивления, рассчитанного по формулам (\ref{eq:group3}) находится в области срабатывания (прямое или обратное направление).
\end{itemize}

\item
Принцип работы линейных органов следующий:
\begin{itemize} 
\item в  качестве  воздействующих  величин  используются  действующие  значения  первой гармоники линейных токов и значения линейных напряжений поляризации рассчитанные по формулам (\ref{eq:group1}); 
\item разрешающими  условиями  для  работы  направленного  органа  являются  превышение  тока значения $0,04 \cdot I_{\textsl{ном}}$ и превышение напряжения значения 1 вольт; 
\item фаза вектора сопротивления, рассчитанного по формулам (\ref{eq:group2}) находится в области срабатывания (прямое или обратное направление).
\end{itemize}

\item
В случае снижения напряжения ниже 0,9 вольт орган направления в течение 100 мс продолжит работу по памяти, и по истечении данного времени произойдет подхват сигнала направления. 

\item
Алгоритм работы органа направления (<<ОН>>) выполнен в соответствии с рисунком \ref{dz:onimg1}. Алгоритм работы логики определения направления, входящий в состав органа направления, выполнен в соответствии с рисунком \ref{dz:onimg2}. В таблице \ref{dz:ontbl1} приведены параметры, необходимые для настройки функции <<ОН>>.

\vspace{3mm}
\begin{figure}[H]
\centering
\includegraphics[width=1\textwidth,height=1\textheight,keepaspectratio]{\fbpath/561.0021 ДЗ/img/on1.pdf}
\captionof{figure}{Функциональная схема логики определения направления}\label{dz:onimg2}
\end{figure}

\small
\begin{longtable}{|>{\centering\arraybackslash}m{5.3cm}|>{\centering\arraybackslash}m{3.3cm}|>{\centering\arraybackslash}m{4.2cm}|>{\centering\arraybackslash}m{1.8cm}|>{\centering\arraybackslash}m{1cm}|}
\caption{Параметры для настройки функции <<ОН>>\hfill\vspace{-0.5\baselineskip}}\label{dz:ontbl1}\\ 
\hline
\rowcolor{gray!30}
Параметр (Параметр на ИЧМ) & Условное обозначение на схеме & Значение/ Диапазон & Единица измерения & Шаг \\ 
\hline
\endfirsthead
\caption*{\hspace{3pt}\emph{Продолжение таблицы \ref{dz:ontbl1}\hfill\vspace{-0.5\baselineskip}}} \\ % сделано по ГОСТ 2.105 п.6.8.7
\hline
\rowcolor{gray!30}
Параметр (Параметр на ИЧМ) & Условное обозначение на схеме & Значение/ Диапазон & Единица измерения & Шаг \\ 
%\hline
\endhead
%\hline
\endfoot
%\hline
\endlastfoot
%==+t1*DISRDIR1|LVDIS
\centering Угол направленности во II квадранте (Фнапр\_II) & \centering φ2 & \centering 90,0 ... 130,0 & \centering град. & \centering \arraybackslash 0,1 \\
\hline
\centering Угол направленности в IV квадранте (Фнапр\_IV) & \centering φ1 & \centering -40,0 ... 0,0 & \centering град. & \centering \arraybackslash 0,1 \\
\hline
%===t1
\end{longtable}
\normalsize

\vspace{3mm}
\begin{figure}[H]
\centering
\includegraphics[width=1\textwidth,height=1\textheight,keepaspectratio]{\fbpath/561.0021 ДЗ/img/on2.pdf}
\captionof{figure}{Функциональная схема <<ОН>>}\label{dz:onimg1}
\end{figure}

\end{enumerate}

%%%%%%%%%%%%%%%%%%%%%%%%%%%%%%%%%%%%%%%%%%%%%%%%%%%%%%%%%%%%%%%%%%%%%% Общие пусковые органы ДЗ (ОПО ДЗ) %%%%%%%%%%%%%%%%%%%%%%%%%%%%%%%%%%%%%%%%%%%%%%%%%%
\item\label{sec:opo} Общие пусковые органы ДЗ (ОПО ДЗ)
\begin{enumerate}[label=\arabic{section}.\arabic{subsection}.\arabic{enumi}.\arabic*, labelsep=4pt, leftmargin=0em, itemindent=65pt, parsep=0pt]

\item
Чтобы  избежать  ложного  срабатывания  ступени  дистанционной  защиты  (ДЗ),  осуществляется дополнительный  контроль  срабатывания  пусковых  органов  дистанционной  защиты  несколькими 
способами:
\begin{itemize}
    \item с пуском по току; 
    \item с пуском по току и напряжению; 
    \item с пуском по току или по току и напряжению.
\end{itemize}
\item
Функция становится активна, если введена в работу одна из ступеней ДЗ и в ней выбран контроль от общих пусковых органов (режим  SGF5 -- <<С контролем от ОПО>> ).
\item
Сигнал общего пуска, пуска от измерительного органа ОПО  ДЗ по току формируется, если максимальное действующее значение первой гармоники одного из фазных токов превысило уставку.
\item
Сигнал общего пуска, пуска от измерительного органа ОПО ДЗ по току и напряжению формируется при снижении одного из линейных напряжений ниже значения уставки и превышении одним из фазных токов значения выше уставки.

\item
Алгоритм работы функции ОПО ДЗ (<<ОПО ДЗ>>) выполнен в соответствии с рисунком \ref{dz:opoimg1}. В таблице \ref{dz:opotbl1} приведены параметры, необходимые для настройки функции <<ОПО ДЗ>>.

\vspace{3mm}
\begin{figure}[H]
\centering
\includegraphics[width=0.8\textwidth,height=0.8\textheight,keepaspectratio]{\fbpath/561.0021 ДЗ/img/opo.pdf}
\captionof{figure}{Функциональная схема <<ОПО ДЗ>>}\label{dz:opoimg1}
\end{figure}

\small
\begin{longtable}{|>{\centering\arraybackslash}m{5.3cm}|>{\centering\arraybackslash}m{3.3cm}|>{\centering\arraybackslash}m{4.2cm}|>{\centering\arraybackslash}m{1.8cm}|>{\centering\arraybackslash}m{1cm}|}
\caption{Параметры для настройки функции <<ОПО ДЗ>>\hfill\vspace{-0.5\baselineskip}}\label{dz:opotbl1}\\ 
\hline
\rowcolor{gray!30}
Параметр (Параметр на ИЧМ) & Условное обозначение на схеме & Значение/ Диапазон & Единица измерения & Шаг \\ 
\hline
\endfirsthead
\caption*{\hspace{3pt}\emph{Продолжение таблицы \ref{dz:opotbl1}\hfill\vspace{-0.5\baselineskip}}} \\ % сделано по ГОСТ 2.105 п.6.8.7
\hline
\rowcolor{gray!30}
Параметр (Параметр на ИЧМ) & Условное обозначение на схеме & Значение/ Диапазон & Единица измерения & Шаг \\ 
%\hline
\endhead
%\hline
\endfoot
%\hline
\endlastfoot
%==+t1*DISPGSU1|LVDIS
\centering Ток срабатывания ПО ДЗ по току (Iср) & \centering I\verb|>>| & \centering 0,10 ... 5,00 & \centering о.е. & \centering \arraybackslash 0,01 \\
\hline
\centering Ток срабатывания ПО ДЗ по току и напряжению (IсрU) & \centering I> & \centering 0,10 ... 5,00 & \centering о.е. & \centering \arraybackslash 0,01 \\
\hline
\centering Напряжение срабатывания ПО по току и напряжению (Uср) & \centering U< & \centering 2,0 ... 100,0 & \centering В & \centering \arraybackslash 0,1 \\
\hline
\centering Режим контроля от БНН (Реж\_БНН) & \centering SGF1 & \centering 0 = Не предусмотрено\\1 = Предусмотрено & \centering -- & \centering \arraybackslash -- \\
\hline
%===t1
\end{longtable}
\normalsize
\end{enumerate}
%%%%%%%%%%%%%%%%%%%%%%%%%%%%%%%%%%%%%%%%%%%%%%%%%%%%%%%%%%%%%%%%%%%%%% КОНЕЦ Общие пусковые органы ДЗ (ОПО ДЗ) %%%%%%%%%%%%%%%%%%%%%%%%%%%%%%%%%%%%%%%%%%%%%%%%%%

%%%%%%%%%%%%%%%%%%%%%%%%%%%%%%%%%%%%%%%%%%%%%%%%%%%%%%%%%%%%%%%%%%%%%%  Определение вида повреждения (ОВП) %%%%%%%%%%%%%%%%%%%%%%%%%%%%%%%%%%%%%%%%%%%%%%%%%%
\item\label{sec:ovp} Определение вида повреждения (ОВП)
\begin{enumerate}[label=\arabic{section}.\arabic{subsection}.\arabic{enumi}.\arabic*, labelsep=4pt, leftmargin=0em, itemindent=65pt, parsep=0pt]

\item    
Данный ИО определяет наличие замыкания на землю и формирует выходной сигнал, указывающий на вид короткого замыкания. 
\item 
Функция становится активной, если введена одна из ступеней ДЗ <<фаза-земля>>. 
\item 
Сигнал о наличии замыкания на землю формируется при выполнении одного из условий:
\begin{itemize} 
    \item отношения  тока  нулевой  последовательности  к  току  прямой  последовательности  превысило уставку по истечению нерегулируемой выдержки времени 100 мс; 
    \item отсутствие сигнала от БННш <<БННш: Срабатывание>> при введенной в работу уставке SGF1; 
    \item действующее значение первой гармоники напряжения НП превысило уставку; 
    \item минимальное действующее значение первой гармоники одного из линейных напряжений ниже уставки. 
\end{itemize}

\item 
Отношение  тока  нулевой  последовательности  к  току  прямой  последовательности  ($3I_0/I_1$) определяет вид КЗ:
\begin{itemize}  
    \item при не срабатывании ПО -- междуфазное КЗ; 
    \item при срабатывании ПО -- двойное КЗ на землю. 
\end{itemize}

\item 
Для ускорения действия функции служит измерительный орган, контролирующий напряжение $3U_0$. При возникновении замыкания на землю в сети с изолированной нейтралью появляется напряжение $3U_0$, 
а также возрастают фазные напряжения неповрежденных фаз, что может вызвать пробой другой фазы на землю и появление двойного замыкания на землю. При срабатывании ИО напряжения $3U_0$ выше уставки 
через  250  мс  происходит  шунтирование  временной  задержки  ИО  тока  в  100  мс,  что  способствует ускорению  срабатывания  функции.  При  этом  условии  осуществляется  контроль исправности  цепей напряжения.
\item 
Для ускорения действия функции так же служит измерительный орган, контролирующий линейные напряжения.  При  снижении  одного  из  линейных  напряжений  ниже  уставки,  шунтируется  выдержка 
времени  ИО  тока  100  мс,  тем  самым  ускоряя  действие  функции.  При  этом  условии  осуществляется контроль исправности цепей напряжения.
\item 
При срабатывании одного из пороговых элементов формируется пусковой сигнал функции <<ОВП: Пуск>>.
\item 
Блокировка  пускового  сигнала  функции  осуществляется  при  наличии введенной  в  работу  уставке контроля от БНН и активного сигнала срабатывания БННш. 

\item
Алгоритм работы функции ОВП (<<ОВП>>) выполнен в соответствии с рисунком \ref{dz:ovpimg1}. В таблице \ref{dz:ovptbl1} приведены параметры, необходимые для настройки функции <<ОВП>>.

\vspace{3mm}
\begin{figure}[H]
\centering
\includegraphics[width=0.9\textwidth,height=0.9\textheight,keepaspectratio]{\fbpath/561.0021 ДЗ/img/ovp.pdf}
\captionof{figure}{Функциональная схема <<ОВП>>}\label{dz:ovpimg1}
\end{figure}

\small
\begin{longtable}{|>{\centering\arraybackslash}m{5.3cm}|>{\centering\arraybackslash}m{3.3cm}|>{\centering\arraybackslash}m{4.2cm}|>{\centering\arraybackslash}m{1.8cm}|>{\centering\arraybackslash}m{1cm}|}
\caption{Параметры для настройки функции <<ОВП>>\hfill\vspace{-0.5\baselineskip}}\label{dz:ovptbl1}\\ 
\hline
\rowcolor{gray!30}
Параметр (Параметр на ИЧМ) & Условное обозначение на схеме & Значение/ Диапазон & Единица измерения & Шаг \\ 
\hline
\endfirsthead
\caption*{\hspace{3pt}\emph{Продолжение таблицы \ref{dz:ovptbl1}\hfill\vspace{-0.5\baselineskip}}} \\ % сделано по ГОСТ 2.105 п.6.8.7
\hline
\rowcolor{gray!30}
Параметр (Параметр на ИЧМ) & Условное обозначение на схеме & Значение/ Диапазон & Единица измерения & Шаг \\ 
%\hline
\endhead
%\hline
\endfoot
%\hline
\endlastfoot
%==+t1*FLOPGSU1|LVDIS
\centering Коэффициент отношения тока нулевой последовательности к току прямой последовательности (Kдел.) & \centering 3I0/I1> & \centering 0,10 ... 2,00 & \centering о.е. & \centering \arraybackslash 0,01 \\
\hline
\centering Уставка по напряжению 3U0 (3U0ср) & \centering 3U0> & \centering 2,0 ... 100,0 & \centering В & \centering \arraybackslash 0,1 \\
\hline
\centering Напряжение срабатывания (Uср) & \centering U< & \centering 2,0 ... 100,0 & \centering В & \centering \arraybackslash 0,1 \\
\hline
\centering Режим контроля от БНН (Реж\_БНН) & \centering SGF1 & \centering 0 = Не предусмотрено\\1 = Предусмотрено & \centering -- & \centering \arraybackslash -- \\
\hline
%===t1
\end{longtable}
\normalsize
\end{enumerate}

%%%%%%%%%%%%%%%%%%%%%%%%%%%%%%%%%%%%%%%%%%%%%%%%%%%%%%%%%%%%%%%%%%%%%% КОНЕЦ  Определение вида повреждения (ОВП) %%%%%%%%%%%%%%%%%%%%%%%%%%%%%%%%%%%%%%%%%%%%%%%%%%

%%%%%%%%%%%%%%%%%%%%%%%%%%%%%%%%%%%%%%%%%%%%%%%%%%%%%%%%%%%%%%%%%%%%%%  Автоматическое ускорение ДЗ (АУ ДЗ)  %%%%%%%%%%%%%%%%%%%%%%%%%%%%%%%%%%%%%%%%%%%%%%%%%%
\item\label{sec:audz} Автоматическое ускорение ДЗ (АУ ДЗ)
\begin{enumerate}[label=\arabic{section}.\arabic{subsection}.\arabic{enumi}.\arabic*, labelsep=4pt, leftmargin=0em, itemindent=65pt, parsep=0pt]

\item    
В  устройстве предусмотрен режим автоматического ускорения ступеней ДЗ.  Выбор ускоряемых ступеней защиты осуществляется из соображений надежного охвата, защищаемого энергообъекта.
\item 
Данный  измерительный орган предусматривает  возможность автоматического  ускорения  пуска  от второй или третьей ступеней ДЗ. 
\item 
Если  уставка  SGF1  выставлена  в  положение  <<2  ступень>>,  то  предусмотрена  возможность автоматического ускорения ДЗ от второй ступени ДЗ при наличии одного из пусковых условий второй ступени <<фаза-фаза>> или 2-й ступени <<фаза-земля>>. 
\item 
Если  уставка  SGF1  выставлена  в  положение  <<3  ступень>>,  то  предусмотрена  возможность автоматического ускорения третьей ступени ДЗ <<фаза-фаза>> при наличии сигнала от пусковых условий данной ступени. Пусковые условия АУ ДЗ формируются без выдержки времени. 

\item
Алгоритм работы функции АУ ДЗ (<<АУ ДЗ>>) выполнен в соответствии с рисунком \ref{dz:audzimg1}. В таблице \ref{dz:audztbl1} приведены параметры, необходимые для настройки функции <<АУ ДЗ>>.

\vspace{3mm}
\begin{figure}[H]
\centering
\includegraphics[width=0.8\textwidth,height=0.8\textheight,keepaspectratio]{\fbpath/561.0021 ДЗ/img/audz.pdf}
\captionof{figure}{Функциональная схема <<АУ ДЗ>>}\label{dz:audzimg1}
\end{figure}

\small
\begin{longtable}{|>{\centering\arraybackslash}m{5.3cm}|>{\centering\arraybackslash}m{3.3cm}|>{\centering\arraybackslash}m{4.2cm}|>{\centering\arraybackslash}m{1.8cm}|>{\centering\arraybackslash}m{1cm}|}
\caption{Параметры для настройки функции <<АУ ДЗ>>\hfill\vspace{-0.5\baselineskip}}\label{dz:audztbl1}\\ 
\hline
\rowcolor{gray!30}
Параметр (Параметр на ИЧМ) & Условное обозначение на схеме & Значение/ Диапазон & Единица измерения & Шаг \\ 
\hline
\endfirsthead
\caption*{\hspace{3pt}\emph{Продолжение таблицы \ref{dz:audztbl1}\hfill\vspace{-0.5\baselineskip}}} \\ % сделано по ГОСТ 2.105 п.6.8.7
\hline
\rowcolor{gray!30}
Параметр (Параметр на ИЧМ) & Условное обозначение на схеме & Значение/ Диапазон & Единица измерения & Шаг \\ 
%\hline
\endhead
%\hline
\endfoot
%\hline
\endlastfoot
%==+t1*AUAPSSU1|LVDIS
\centering Выбор ускоряемой ступени (Выбор\_ст) & \centering SGF1 & \centering 1 = 2 ступень\\2 = 3 ступень & \centering -- & \centering \arraybackslash -- \\
\hline
%===t1
\end{longtable}
\normalsize

\end{enumerate}
%%%%%%%%%%%%%%%%%%%%%%%%%%%%%%%%%%%%%%%%%%%%%%%%%%%%%%%%%%%%%%%%%%%%%% КОНЕЦ  Автоматическое ускорение ДЗ (АУ ДЗ) %%%%%%%%%%%%%%%%%%%%%%%%%%%%%%%%%%%%%%%%%%%%%%%%%%


%%%%%%%%%%%%%%%%%%%%%%%%%%%%%%%%%%%%%%%%%%%%%%%%%%%%%%%%%%%%%%%%%%%%%%  Оперативное ускорение ДЗ (ОУ ДЗ)  %%%%%%%%%%%%%%%%%%%%%%%%%%%%%%%%%%%%%%%%%%%%%%%%%%
\item\label{sec:oudz} Оперативное ускорение ДЗ (ОУ ДЗ)
\begin{enumerate}[label=\arabic{section}.\arabic{subsection}.\arabic{enumi}.\arabic*, labelsep=4pt, leftmargin=0em, itemindent=65pt, parsep=0pt]

\item   
Оперативное ускорение ДЗ может вводиться при:
\begin{itemize}
    \item выводе основной защиты оборудования; 
    \item опробовании оборудования рабочим напряжением; 
    \item операциях с разъединителями на противоположном конце ВЛ. 
\end{itemize}

\item 
В  данном  алгоритме  активном состоянии входа <<Ввод  ОУ  ДЗ>>,  уставкой  SGF1  предусмотрена возможность выбора оперативного ускорения от пуска одной из ступеней ДЗ. 
\item
Если  уставка  SGF1  выставлена  в  положение  <<2  ступень>>,  то  предусмотрена  возможность автоматического ускорения ДЗ от второй ступени ДЗ при наличии одного из пусковых условий второй ступени <<фаза-фаза>> или 2-й ступени <<фаза-земля>>.  
\item
Если  уставка  SGF1  выставлена  в  положение  <<3  ступень>>,  то  предусмотрена  возможность автоматического ускорения ДЗ от 3-й ступени ДЗ <<фаза-фаза>> при наличии сигнала от пусковых условий данной ступени.
\item
Срабатывание оперативного ускорения от пуска ступеней ДЗ происходит по истечению выдержки времени таймера Т1. 

\item
Алгоритм работы функции ОУ ДЗ (<<ОУ ДЗ>>) выполнен в соответствии с рисунком \ref{dz:oudzimg1}. В таблице \ref{dz:oudztbl1} приведены параметры, необходимые для настройки функции <<ОУ ДЗ>>.

\vspace{3mm}
\begin{figure}[H]
\centering
\includegraphics[width=0.8\textwidth,height=0.8\textheight,keepaspectratio]{\fbpath/561.0021 ДЗ/img/oudz.pdf}
\captionof{figure}{Функциональная схема <<ОУ ДЗ>>}\label{dz:oudzimg1}
\end{figure}

\small
\begin{longtable}{|>{\centering\arraybackslash}m{5.3cm}|>{\centering\arraybackslash}m{3.3cm}|>{\centering\arraybackslash}m{4.2cm}|>{\centering\arraybackslash}m{1.8cm}|>{\centering\arraybackslash}m{1cm}|}
\caption{Параметры для настройки функции <<ОУ ДЗ>>\hfill\vspace{-0.5\baselineskip}}\label{dz:oudztbl1}\\ 
\hline
\rowcolor{gray!30}
Параметр (Параметр на ИЧМ) & Условное обозначение на схеме & Значение/ Диапазон & Единица измерения & Шаг \\ 
\hline
\endfirsthead
\caption*{\hspace{3pt}\emph{Продолжение таблицы \ref{dz:oudztbl1}\hfill\vspace{-0.5\baselineskip}}} \\ % сделано по ГОСТ 2.105 п.6.8.7
\hline
\rowcolor{gray!30}
Параметр (Параметр на ИЧМ) & Условное обозначение на схеме & Значение/ Диапазон & Единица измерения & Шаг \\ 
%\hline
\endhead
%\hline
\endfoot
%\hline
\endlastfoot
%==+t1*OPAPSSU1|LVDIS
\centering Выбор оперативно ускоряемой ступени (Выбор\_ст\_ОУ) & \centering SGF1 & \centering 1 = 2 ступень\\2 = 3 ступень & \centering -- & \centering \arraybackslash -- \\
\hline
\centering Выдержка времени срабатывания оперативного ускорения (T\_ОУ) & \centering T1 & \centering 0,00 ... 3,00 & \centering с & \centering \arraybackslash 0,01 \\
\hline
%===t1
\end{longtable}
\normalsize

\end{enumerate}
%%%%%%%%%%%%%%%%%%%%%%%%%%%%%%%%%%%%%%%%%%%%%%%%%%%%%%%%%%%%%%%%%%%%%% КОНЕЦ  Оперативное ускорение ДЗ (ОУ ДЗ) %%%%%%%%%%%%%%%%%%%%%%%%%%%%%%%%%%%%%%%%%%%%%%%%%%

%%%%%%%%%%%%%%%%%%%%%%%%%%%%%%%%%%%%%%%%%%%%%%%%%%%%%%%%%%%%%%%%%%%%%%  Внешняя логика функции ДЗ   %%%%%%%%%%%%%%%%%%%%%%%%%%%%%%%%%%%%%%%%%%%%%%%%%%
\item\label{sec:vldz} Внешняя логика функции ДЗ
\begin{enumerate}[label=\arabic{section}.\arabic{subsection}.\arabic{enumi}.\arabic*, labelsep=4pt, leftmargin=0em, itemindent=65pt, parsep=0pt]

\item
Для функции предусмотрен общий выходной сигнал «ДЗ: Пуск», который формируется при наличии пусковых сигналов от одной из ступеней ДЗ. 

\item
Алгоритм работы  внешней логики ДЗ выполнен в соответствии с рисунком \ref{dz:vlimg1}.

\vspace{3mm}
\begin{figure}[H]
\centering
\includegraphics[width=0.6\textwidth,height=0.6\textheight,keepaspectratio]{\fbpath/561.0021 ДЗ/img/vldz.pdf}
\captionof{figure}{Функциональная схема внешней логики ДЗ}\label{dz:vlimg1}
\end{figure}

\end{enumerate}
%%%%%%%%%%%%%%%%%%%%%%%%%%%%%%%%%%%%%%%%%%%%%%%%%%%%%%%%%%%%%%%%%%%%%% КОНЕЦ  Внешняя логика функции ДЗ %%%%%%%%%%%%%%%%%%%%%%%%%%%%%%%%%%%%%%%%%%%%%%%%%%


\end{enumerate}
